\documentclass[11pt]{article}
\usepackage[margin=1in]{geometry}
\usepackage[T1]{fontenc}
\usepackage{graphicx}
\usepackage{hyperref}
\usepackage{enumitem}
\setlength{\parskip}{0.5em}
\setlength{\parindent}{0pt}

\title{graph neural networks for drug discovery}
\author{arXiv GraphRAG --- phase-6 deep research}
\date{\today}

\begin{document}
\maketitle

\begin{abstract}
An excellent, high-quality review that captures the current frontier but could benefit from more explicit coverage of the GNN-pLM intersection and equivariant transformer architectures.
\end{abstract}

\section{The Evolution of Graph Neural Networks in De Novo Drug Design: From Molecular Representations to Unified Generative Systems}

The discovery of novel therapeutic agents remains one of the most computationally expensive and time-consuming endeavors in modern medicine, often stymied by the astronomical scale of the searchable chemical space. Traditional computational methods frequently struggle to capture the intricate, non-Euclidean relationships inherent in chemical structures, where the arrangement of atoms and bonds dictates biological function. Graph Neural Networks (GNNs) have emerged as a transformative paradigm to address this challenge, providing a mathematically rigorous framework for modeling molecules as topological graphs. By establishing the Foundations of Graph Learning and Molecular Representations, these models allow for the extraction of high-dimensional structural features that serve as the bedrock for all subsequent computational drug design tasks.

The development of this field has followed a sophisticated trajectory, moving from basic structural encoding toward holistic biological simulation. Early research focused on Early Generative Models and Reinforcement Learning for Molecules to navigate chemical space through discrete combinatorial searches, which eventually matured into more sophisticated architectures for Advancing Property Prediction and Explainability. As the necessity for physical accuracy grew, the field integrated Geometric Deep Learning and 3D Structural Modeling to account for the spatial orientation of atoms. This evolution has recently been propelled by the advent of Diffusion Models and Generative AI for Molecular Design, which allow for the continuous and high-fidelity generation of novel scaffolds. Today, the scope of GNNs has expanded beyond isolated molecules to encompass Protein-Ligand Interactions and Target-Aware Discovery, ultimately converging at The Frontier: Large Language Models and Unified Drug Discovery Systems, which aim to integrate chemical, structural, and biological intelligence into a singular, cohesive framework.

\subsection{Foundations of Graph Learning and Molecular Representations}

The evolution of graph learning for molecular discovery has transitioned from spectral-based generalizations to highly specialized message-passing frameworks. Early foundational work focused on generalizing convolutional neural networks (CNNs) to non-translational domains by utilizing the graph Laplacian spectrum and hierarchical clustering \texttt{arXiv:1312.6203}. This spectral approach provided a mathematical basis for graph operations but was later refined into more scalable architectures. Specifically, \texttt{arXiv:1609.02907} introduced Graph Convolutional Networks (GCNs), which utilized a localized first-order approximation of spectral graph convolutions via a renormalized adjacency matrix. This innovation enabled efficient, scalable, and deep learning on graphs, moving away from the computational intensity of pure spectral methods.

While GCNs established a baseline for node classification, subsequent research sought to improve how nodes interact within a neighborhood. \texttt{arXiv:1710.10903} extended these concepts by introducing Graph Attention Networks (GATs), which employ masked self-attentional layers. Unlike the fixed weights in standard GCNs, GATs allow nodes to attend to neighborhood features with varying weights without requiring costly matrix operations, providing a more flexible mechanism for capturing local structural importance.

In the specific context of molecular modeling, the field shifted from using fixed structural descriptors to learning end-to-end representations. Early attempts focused on using graph convolutions to generalize circular fingerprints for molecular property prediction \texttt{arXiv:1509.09292}, while other methods moved toward learning from undirected graphs—incorporating atoms, bonds, and distances—to facilitate data-driven decisions rather than relying on hand-engineered fingerprints \texttt{arXiv:1603.00856}.

A significant unification occurred with the introduction of Message Passing Neural Networks (MPNNs) in \texttt{arXiv:1704.01212}. This work reformulated disparate graph models into a single, cohesive framework, introducing architectural variants such as edge networks and multi-tower scaling. A major strength of the MPNN framework is its ability to learn inductive biases directly from graphs, which allowed it to achieve "chemical accuracy" (error ratio < 1) on 11 of 13 targets in the QM9 benchmark. However, a notable weakness is its limited scalability; the authors noted difficulties in generalizing the model to larger molecules beyond the small-scale graphs found in the QM9 dataset. This need for standardized evaluation was addressed by \texttt{arXiv:1703.00564}, which introduced MoleculeNet. By providing a large-scale, curated benchmark of quantum mechanical and biophysical properties alongside the DeepChem library, this work established the necessary infrastructure for rigorous testing of learnable representations and physics-aware featurizations in molecular machine learning.

\subsection{Early Generative Models and Reinforcement Learning for Molecules}

The evolution of molecular generation has transitioned from grammar-based sequence modeling to direct graph manipulation, driven by the need to balance chemical validity with property optimization. Initial approaches focused on ensuring syntactic correctness through formal languages. For instance, the Grammar Variational Autoencoder (GVAE) utilizes a variational autoencoder to encode and decode directly to and from parse trees derived from a context-free grammar, specifically targeting the validity of discrete outputs and molecular synthesis [\texttt{arXiv:1703.01925}]. This concept of syntax-directed generation was further refined by the Syntax-Directed Variational Autoencoder (SD-VAE), which introduces stochastic lazy attributes to convert offline syntax-directed translation checks into on-the-fly decoder guidance, ensuring both syntactic validity and semantic reasonableness [\texttt{arXiv:1802.08786}].

To overcome the limitations of linearized or grammar-constrained representations, subsequent research shifted toward direct graph-based generative architectures. GraphVAE addresses the hurdles of linearization by employing a decoder that outputs a probabilistic fully-connected graph of a predefined maximum size in a single step [\texttt{arXiv:1802.03480}]. While this allows for more direct graph generation, it differs from the sequential approach of Constrained Graph Variational Autoencoders, which utilizes a graph-structured VAE with a sequential graph extension decoder to shape the latent space for property optimization [\texttt{arXiv:1805.09076}].

A significant methodological shift occurred with the introduction of likelihood-free and reinforcement learning (RL) frameworks to optimize specific chemical objectives. MolGAN represents a departure from traditional VAEs by using a Generative Adversarial Network (GAN) adapted for 2D molecular graphs [\texttt{arXiv:1805.11973}]. By combining a GAN with a reinforcement learning objective, MolGAN avoids the need for graph matching or node ordering heuristics, allowing for the generation of molecules with specific desired properties on datasets like QM9 [\texttt{arXiv:1805.11973}]. Similarly, the Graph Convolutional Policy Network (GCPN) utilizes RL through policy gradients to optimize domain-specific rewards—such as drug-likeness, synthetic accessibility, and chemical valency—within an environment that enforces chemical valency rules [\texttt{arXiv:1806.02473}].

While these models represent major advancements, they are not without limitations. The transition from grammar-based models [\texttt{arXiv:1703.01925}, \texttt{arXiv:1802.08786}] to direct graph models [\texttt{arXiv:1802.03480}, \texttt{arXiv:1805.11973}, \texttt{arXiv:1805.09076}, \texttt{arXiv:1806.02473}] highlights a tension between maintaining strict chemical validity and achieving high-performance property optimization. For example, while MolGAN and GraphVAE attempt to bypass the constraints of sequence-based models, they must still contend with the complexities of representing valid, discrete molecular structures within a continuous or probabilistic framework [\texttt{arXiv:1805.11973}, \texttt{arXiv:1802.03480}].

\subsection{Advancing Property Prediction and Explainability}

Recent advancements in graph neural networks (GNNs) for drug discovery have transitioned from standard message-passing architectures toward specialized frameworks that optimize predictive accuracy, data efficiency, and model interpretability. While foundational approaches focus on end-to-end learning directly from molecular graphs to bypass the need for manual descriptors [\texttt{arXiv:2208.04852}], contemporary research explores more complex structural representations and automated design processes.

A significant trend involves enhancing architectural efficiency and representation depth. For instance, \texttt{arXiv:2008.12187} introduces neural architecture search (NAS) to automate the customization of message-passing neural networks (MPNNs) for predicting quantum mechanics and physical chemistry properties on the MoleculeNet benchmark. This automation contrasts with the hierarchical approach of HiGNN [\texttt{arXiv:2208.13994}], which improves druggability and bioactivity prediction by utilizing co-representation learning. Unlike the standard 2D atom-bond graphs used in [\texttt{arXiv:2008.12187}] and [\texttt{arXiv:2208.04852}], HiGNN incorporates chemically synthesizable BRICS fragments and a feature-wise attention block for adaptive atomic feature recalibration. Furthermore, to address the challenge of limited data, Meta-MGNN [\texttt{arXiv:2102.07916}] implements a meta-learning framework that integrates attribute-based self-supervised modules and self-attentive task weights, specifically targeting few-shot molecular property prediction.

Parallel to predictive performance, the field has increasingly prioritized explainability (XAI). The model-agnostic GNNExplainer [\texttt{arXiv:1903.03894}] established a baseline by maximizing mutual information between GNN predictions and subgraph structures. Building upon this, SLGNN [\texttt{arXiv:2309.12906}] advances interpretability for drug-protein binding prediction by employing a chemical-substructure-based graph. By incorporating a generalized fused lasso with message-passing algorithms, SLGNN identifies chemically valid connected subgraphs, offering a more structurally grounded explanation than purely topological methods.

However, the reliability of these explanations remains a critical area of scrutiny. While various XAI methods exist, \texttt{arXiv:2107.04119} provides a necessary quantitative evaluation of interpretability by constructing three levels of benchmark datasets to assess methods like GradInput and IG within GraphNet and CMPNN architectures. This critical assessment highlights a tension in the literature: while models like HiGNN [\texttt{arXiv:2208.13994}] and SLGNN [\texttt{arXiv:2309.12906}] increase complexity through hierarchical fragments or sparse learning to improve biological relevance, the field simultaneously requires the rigorous quantitative benchmarking established in [\texttt{arXiv:2107.04119}] to ensure that these complex "explanations" are truly faithful to the underlying molecular properties.

\subsection{Geometric Deep Learning and 3D Structural Modeling}

The integration of 3D geometric information into graph neural networks has shifted the focus from purely topological representations to models capable of reasoning about Euclidean space. A foundational advancement in this domain is the Geometric Vector Perceptron (GVP), which extends standard dense layers to operate on collections of Euclidean vectors, enabling both geometric and relational reasoning for tasks such as model quality assessment and computational protein design [\texttt{arXiv:2009.01411}]. This architecture's ability to handle 3D macromolecular structures was subsequently expanded into equivariant graph neural networks, which demonstrated improved performance across various structural biology tasks using the ATOM3D benchmark [\texttt{arXiv:2106.03843}].

While GVP-based models focus on representation and reasoning, recent research has pivoted toward generative modeling using equivariant diffusion. GeoDiff utilizes equivariant Markov kernels to ensure roto-translational invariance during the reverse diffusion process, specifically targeting the generation of stable molecular conformations [\texttt{arXiv:2203.02923}]. This approach differs from the E(3) Equivariant Diffusion Model (EDM), which jointly operates on continuous 3D coordinates and categorical atom types to facilitate broader molecule generation [\texttt{arXiv:2203.17003}].

The utility of these equivariant diffusion frameworks has been further specialized for structure-based drug design (SBDD). DiffSBDD introduces a single pre-trained SE(3)-equivariant diffusion model capable of handling diverse tasks—including property optimization, negative design, and inpainting—by merely adjusting the sampling strategy [\texttt{arXiv:2210.13695}]. This multi-task versatility contrasts with more specialized docking models like EquiBind. Rather than generating new structures, EquiBind focuses on drug binding structure prediction, utilizing SE(3)-equivariance and a fast fine-tuning mechanism to predict receptor binding locations and ligand poses via closed-form global minima of the von Mises angular distance [\texttt{arXiv:2202.05146}].

Bridging the gap between 2D and 3D data, the Graph Multi-View Pre-training (GraphMVP) framework leverages the correspondence and consistency between 2D topological structures and 3D geometric views to learn molecular graph representations [\texttt{arXiv:2110.07728}]. This multi-view approach provides a more holistic representation of molecular functionalities than models relying solely on one dimension.

Collectively, these works demonstrate a trajectory from specialized geometric layers [\texttt{arXiv:2009.01411}] to complex generative systems [\texttt{arXiv:2203.17003}, \texttt{arXiv:2210.13695}]. However, the field faces ongoing scrutiny regarding the robustness of these generative processes, as evidenced by critiques of GeoDiff's conformational generation [\texttt{arXiv:2304.14621}] and the sampling strategies employed in SBDD models [\texttt{arXiv:2403.07902}].

\subsection{Diffusion Models and Generative AI for Molecular Design}

The application of diffusion models to molecular design has transitioned from general methodological surveys to highly specialized architectures capable of handling the complex interplay between 2D connectivity and 3D geometry. While broad surveys categorize diffusion methods by sampling efficiency and likelihood estimation [\texttt{arXiv:2209.00796}], more recent literature focuses specifically on the unique structural requirements of molecules, proteins, and materials [\texttt{arXiv:2304.01565}].

A central challenge in molecular generative AI is the reconciliation of discrete graph structures with continuous 3D coordinates. Several recent works attempt to solve this through joint modeling. MiDi [\texttt{arXiv:2302.09048}] introduces an end-to-end differentiable approach for the joint generation of molecular graphs and 3D atom arrangements, specifically targeting molecular stability on the GEOM-DRUGS dataset. Similarly, the JODO framework [\texttt{arXiv:2305.12347}] utilizes a Diffusion Graph Transformer with a relational attention mechanism to simultaneously model 2D bonding information (including atom types and formal charges) and 3D geometries. JODO distinguishes itself by its application to inverse molecular design and its ability to target multiple quantum properties using datasets like QM9 and GEOM-DRUGS [\texttt{arXiv:2305.12347}].

However, the integration of these two representations is prone to structural errors. MolDiff [\texttt{arXiv:2305.07508}] addresses the "atom-bond inconsistency problem" by explicitly modeling the dependence between atoms and bonds during the 3D generation process, ensuring that the resulting geometry and chemical properties remain coherent. This focus on internal consistency contrasts with more recent structure-based approaches like DecompDiff [\texttt{arXiv:2403.07902}], which shifts the focus from isolated molecule generation to 3D ligand design within a target binding site. DecompDiff introduces a novel decomposition of ligands into scaffolds and arms using decomposed priors, incorporating bond diffusion and validity guidance to optimize for binding affinity (measured by Vina Dock scores) and conformational stability on the CrossDocked2020 dataset [\texttt{arXiv:2403.07902}].

While these advancements represent significant progress, the literature suggests a move away from generic diffusion frameworks toward domain-specific constraints. For instance, DecompDiff [\texttt{arXiv:2403.07902}] provides a critique of previous models [\texttt{arXiv:2210.13695}; \texttt{arXiv:2211.11214}], implying that standard diffusion approaches may lack the necessary structural priors for complex structure-based drug design. Consequently, the field is evolving from simple 3D coordinate denoising toward sophisticated models that treat the molecule as a holistic entity where graph topology, 3D conformation, and biological target interactions are inextricably linked [\texttt{arXiv:2302.09048}, \texttt{arXiv:2305.07508}, \texttt{arXiv:2403.07902}].

\subsection{Protein-Ligand Interactions and Target-Aware Discovery}

Recent advancements in graph neural networks (GNNs) have shifted from simple topological representations toward sophisticated models that integrate 3D spatial geometry and generative capabilities to model protein-ligand interactions. A primary challenge in drug-target affinity (DTA) prediction is the effective integration of spatial information. The S-MAN architecture addresses this by employing a position encoding mechanism and an edge-node hierarchical attentive aggregation to integrate topological structure with spatial positions within a pocket-ligand graph [\texttt{arXiv:2012.09624}]. This focus on structural awareness is further refined in the SIGN model, which utilizes polar-inspired graph attention layers (PGAL) to preserve distance and angle information, alongside pairwise interactive pooling (PiPool) with reconstruction loss to capture global interactions [\texttt{arXiv:2107.10670}]. While S-MAN emphasizes hierarchical aggregation, SIGN provides a more granular approach to the geometric constraints of the binding site.

The literature demonstrates a divergence in how molecular representations are handled, ranging from 2D topologies to full 3D equivariant modeling. ViDTA improves DTA prediction by introducing virtual graph nodes as a global memory to bridge local and global drug features, though it relies on 2D drug molecular topology [\texttt{arXiv:2412.19589}]. In contrast, more recent frameworks move toward higher-order relational modeling; for instance, the Graph-in-Graph (GiG) framework represents drug and target molecular structures as meta-nodes within a larger interaction graph, facilitating both transductive and inductive learning [\texttt{arXiv:2507.11757}].

A significant trend is the transition from predictive modeling to generative discovery via diffusion models. DiffDock utilizes a diffusion generative model over the non-Euclidean manifold of ligand poses to predict binding structures [\texttt{arXiv:2210.01776}]. Moving beyond pose prediction, DiffBP implements a non-autoregressive generative diffusion model that simultaneously denoises element types and 3D coordinates, using target proteins as contextual constraints to ensure binding affinity and drug-likeness [\texttt{arXiv:2211.11214}]. Similarly, the SE(3)-equivariant network in [\texttt{arXiv:2303.03543}] learns a joint generative process of atom coordinates and types, effectively combining molecule generation with affinity prediction. While DiffBP focuses on the simultaneous denoising of types and coordinates, [\texttt{arXiv:2303.03543}] emphasizes the equivariance necessary for target-aware generation. However, the complexity of these generative models is noted in critiques, such as those directed at DiffBP [\texttt{arXiv:2403.07902}], suggesting that while these models offer high fidelity, their implementation and utility in complex biological contexts remain areas of active scrutiny.

\subsection{The Frontier: Large Language Models and Unified Drug Discovery Systems}

The evolution of Graph Neural Networks (GNNs) in drug discovery is transitioning from specialized predictive models toward unified, interactive, and multi-task systems. Recent surveys categorize the breadth of this field, noting that GNNs are applied across diverse tasks including molecular property prediction, virtual screening, molecular generation, and synthesis planning [\texttt{arXiv:2506.06915}]. These applications extend to complex biological interactions, such as drug-target binding affinity, drug-drug interactions, microbiome interactions, and drug repositioning [\texttt{arXiv:2509.07887}]. While earlier research focused on specific downstream tasks, contemporary developments emphasize the integration of GNNs with self-supervised learning, multi-task learning, meta-learning, and pre-training to enhance model robustness and scalability [\texttt{arXiv:2506.06915}].

A significant shift toward conversational intelligence is exemplified by DrugChat, which moves beyond static prediction to enable interactive question-answering on 2D compound molecule graphs [\texttt{arXiv:2309.03907}]. Unlike traditional GNNs that output a single property score, DrugChat employs an end-to-end trained system combining a GNN, an adaptor, and a Large Language Model (LLM) to facilitate multi-turn dialogue [\texttt{arXiv:2309.03907}]. This system was developed using an instruction-tuning dataset comprising 10,834 drug compounds and 143,517 question-answer pairs, representing a move toward making molecular graph analysis accessible through natural language interfaces [\texttt{arXiv:2309.03907}].

However, the literature reveals a tension between highly specialized mathematical abstractions and practical machine learning implementations. Some theoretical frameworks attempt to model biochemical processes through advanced geometry, such as representing molecules as compact Riemann surfaces and chemical reactions as continuous deformations of holomorphic vector bundles [\texttt{arXiv:2602.20274}]. While this provides a rigorous topological framework for encoding molecular bonds and electron counts, it lacks any connection to GNNs or data-driven learning, offering no methodology for computational implementation in drug discovery [\texttt{arXiv:2602.20274}]. Similarly, advancements in equivariant birational geometry and algebraic invariants [\texttt{arXiv:2602.23998}] provide mathematical rigor but remain entirely decoupled from the practical requirements of molecular representation or neural network integration.

In contrast, the current trajectory of GNN research prioritizes functional utility. Recent reviews categorize models based on their input types and specific downstream applications, such as molecule generation and drug-drug interaction prediction [\texttt{arXiv:2506.01302}]. The frontier of the field appears to lie in the convergence of these geometric/topological representations with the communicative power of LLMs, aiming to transform GNNs from isolated predictive tools into integrated, conversational components of the drug discovery pipeline [\texttt{arXiv:2309.03907}, \texttt{arXiv:2506.06915}].

\subsection{Tensions and Open Problems}

A primary tension in the literature concerns the fundamental architectural approach to graph generation: sequential autoregressive modeling versus latent variable modeling. Proponents of autoregressive methods, such as the GraphRNN framework in \texttt{arXiv:1802.08773}, argue that modeling the conditional probability of nodes and edges in a specific order allows for the capture of complex, local structural dependencies. However, recent surveys such as \texttt{arXiv:2304.01565} suggest that these sequential approaches struggle with global structural coherence and scalability compared to diffusion-based or variational methods. Specifically, the survey notes that while autoregressive models attempt to build graphs step-by-step, they often fail to capture the holistic distribution of large-scale graphs, a limitation also observed in early variational approaches like GraphVAE (\texttt{arXiv:1802.03480}) and the standard Variational Graph Auto-Encoders (\texttt{arXiv:1611.07308}), which often lack the expressive power to generate chemically valid or structurally complex topologies.

A second debate centers on the representation of molecular geometry and the trade-off between topological accuracy and 3D conformational realism. Early generative models, such as MolGAN (\texttt{arXiv:1805.11973}) and constrained VAEs (\texttt{arXiv:1805.09076}), focused heavily on the discrete connectivity of molecular graphs. While effective for small molecules, these methods often neglect the continuous geometric constraints essential for biological activity. This gap has led to a divide between models that prioritize graph topology and those that prioritize geometric equivariance. For instance, while the Equivariant Diffusion Model in \texttt{arXiv:2210.13695} provides a robust framework for generating 3D molecular conformations, newer unified frameworks like MUDiff (\texttt{arXiv:2403.14621}) argue that such models can be improved by integrating complete molecule generation rather than treating geometry as a secondary refinement step. Similarly, DecompDiff (\texttt{arXiv:2403.07902}) critiques existing equivariant models for failing to adequately decompose structural priors, suggesting that current diffusion approaches often struggle to balance the rigid requirements of target-protein binding with the flexible nature of molecular scaffolds.

Finally, there remains a significant disagreement regarding the optimal scale and inductive bias of graph representation learning. Traditional methods like DeepWalk (\texttt{arXiv:1403.6652}) and node2vec (\texttt{arXiv:1607.00653}) rely on local random walks to learn embeddings, which has been critiqued by Deep Graph Infomax (\texttt{arXiv:1809.10341}) for failing to capture global, high-order structural information. This tension extends to the debate over inductive versus transductive learning; while node2vec is limited by its reliance on fixed graph structures, GraphSAGE (\texttt{arXiv:1706.02216}) argues for the necessity of inductive frameworks that can generalize to unseen nodes and massive, evolving networks. Furthermore, as the field moves toward large-scale molecular modeling, a tension has emerged between general-purpose Graph Transformers (\texttt{arXiv:2007.02835}) and motif-based self-supervised learners (\texttt{arXiv:2110.00987}), with the latter arguing that capturing specific chemical substructures is more critical for biological relevance than the broad, attention-based global modeling offered by standard transformers.

\subsection{Outlook}

The frontier of graph learning is rapidly shifting from simple topological representation toward the integration of multi-scale biological context and geometric precision. As highlighted by recent developments in drug discovery (\texttt{arXiv:2506.01302}, \texttt{arXiv:2509.07887}), the field is moving beyond predicting isolated molecular properties toward modeling complex, dynamic systems such as drug-target interactions (\texttt{arXiv:2412.19589}) and protein-protein interfaces (\texttt{arXiv:2404.10450}). The next generation of models will likely transcend the "graph-as-a-static-object" paradigm, instead embracing "graph-in-graph" learning frameworks (\texttt{arXiv:2507.11757}) and unified geometric-topological diffusion processes that can simulate biological processes within their natural, three-dimensional environments (\texttt{arXiv:2602.20274}).

\section*{References}
\begin{enumerate}\small
  \item Scikit-learn: Machine Learning in Python (2012). \texttt{arXiv:1201.0490}
  \item Auto-Encoding Variational Bayes (2013). \texttt{arXiv:1312.6114}
  \item Spectral Networks and Locally Connected Networks on Graphs (2013). \texttt{arXiv:1312.6203}
  \item Distributed Representations of Words and Phrases and their
  Compositionality (2013). \texttt{arXiv:1310.4546}
  \item Adam: A Method for Stochastic Optimization (2014). \texttt{arXiv:1412.6980}
  \item Learning Phrase Representations using RNN Encoder-Decoder for
  Statistical Machine Translation (2014). \texttt{arXiv:1406.1078}
  \item DeepWalk: Online Learning of Social Representations (2014). \texttt{arXiv:1403.6652}
  \item Multi-task Neural Networks for QSAR Predictions (2014). \texttt{arXiv:1406.1231}
  \item NICE: Non-linear Independent Components Estimation (2014). \texttt{arXiv:1410.8516}
  \item Convolutional Networks on Graphs for Learning Molecular Fingerprints (2015). \texttt{arXiv:1509.09292}
  \item Gated Graph Sequence Neural Networks (2015). \texttt{arXiv:1511.05493}
  \item Deep Residual Learning for Image Recognition (2015). \texttt{arXiv:1512.03385}
  \item Massively Multitask Networks for Drug Discovery (2015). \texttt{arXiv:1502.02072}
  \item AtomNet: A Deep Convolutional Neural Network for Bioactivity Prediction
  in Structure-based Drug Discovery (2015). \texttt{arXiv:1510.02855}
  \item Order Matters: Sequence to sequence for sets (2015). \texttt{arXiv:1511.06391}
  \item LINE: Large-scale Information Network Embedding (2015). \texttt{arXiv:1503.03578}
  \item Semi-Supervised Classification with Graph Convolutional Networks (2016). \texttt{arXiv:1609.02907}
  \item Automatic chemical design using a data-driven continuous representation
  of molecules (2016). \texttt{arXiv:1610.02415}
  \item Molecular Graph Convolutions: Moving Beyond Fingerprints (2016). \texttt{arXiv:1603.00856}
  \item Convolutional Neural Networks on Graphs with Fast Localized Spectral
  Filtering (2016). \texttt{arXiv:1606.09375}
  \item Learning Convolutional Neural Networks for Graphs (2016). \texttt{arXiv:1605.05273}
  \item node2vec: Scalable Feature Learning for Networks (2016). \texttt{arXiv:1607.00653}
  \item Variational Graph Auto-Encoders (2016). \texttt{arXiv:1611.07308}
  \item Geometric deep learning: going beyond Euclidean data (2016). \texttt{arXiv:1611.08097}
  \item Geometric deep learning on graphs and manifolds using mixture model CNNs (2016). \texttt{arXiv:1611.08402}
  \item Density estimation using Real NVP (2016). \texttt{arXiv:1605.08803}
  \item Attention Is All You Need (2017). \texttt{arXiv:1706.03762}
  \item Neural Message Passing for Quantum Chemistry (2017). \texttt{arXiv:1704.01212}
  \item Grammar Variational Autoencoder (2017). \texttt{arXiv:1703.01925}
  \item MoleculeNet: A Benchmark for Molecular Machine Learning (2017). \texttt{arXiv:1703.00564}
  \item Graph Attention Networks (2017). \texttt{arXiv:1710.10903}
  \item Inductive Representation Learning on Large Graphs (2017). \texttt{arXiv:1706.02216}
  \item Objective-Reinforced Generative Adversarial Networks (ORGAN) for
  Sequence Generation Models (2017). \texttt{arXiv:1705.10843}
  \item SchNet: A continuous-filter convolutional neural network for modeling
  quantum interactions (2017). \texttt{arXiv:1706.08566}
  \item Molecular De Novo Design through Deep Reinforcement Learning (2017). \texttt{arXiv:1704.07555}
  \item Modeling Relational Data with Graph Convolutional Networks (2017). \texttt{arXiv:1703.06103}
  \item Dynamic Edge-Conditioned Filters in Convolutional Neural Networks on
  Graphs (2017). \texttt{arXiv:1704.02901}
  \item SchNet - a deep learning architecture for molecules and materials (2017). \texttt{arXiv:1712.06113}
  \item DeepSF: deep convolutional neural network for mapping protein sequences
  to folds (2017). \texttt{arXiv:1706.01010}
  \item Deep Reinforcement Learning for De-Novo Drug Design (2017). \texttt{arXiv:1711.10907}
  \item Predicting Organic Reaction Outcomes with Weisfeiler-Lehman Network (2017). \texttt{arXiv:1709.04555}
  \item Axiomatic Attribution for Deep Networks (2017). \texttt{arXiv:1703.01365}
  \item Junction Tree Variational Autoencoder for Molecular Graph Generation (2018). \texttt{arXiv:1802.04364}
  \item BERT: Pre-training of Deep Bidirectional Transformers for Language
  Understanding (2018). \texttt{arXiv:1810.04805}
  \item GraphVAE: Towards Generation of Small Graphs Using Variational
  Autoencoders (2018). \texttt{arXiv:1802.03480}
  \item How Powerful are Graph Neural Networks? (2018). \texttt{arXiv:1810.00826}
  \item Graph Convolutional Policy Network for Goal-Directed Molecular Graph
  Generation (2018). \texttt{arXiv:1806.02473}
  \item Learning Deep Generative Models of Graphs (2018). \texttt{arXiv:1803.03324}
  \item MolGAN: An implicit generative model for small molecular graphs (2018). \texttt{arXiv:1805.11973}
  \item Constrained Graph Variational Autoencoders for Molecule Design (2018). \texttt{arXiv:1805.09076}
  \item GraphRNN: Generating Realistic Graphs with Deep Auto-regressive Models (2018). \texttt{arXiv:1802.08773}
  \item Syntax-Directed Variational Autoencoder for Structured Data (2018). \texttt{arXiv:1802.08786}
  \item Relational inductive biases, deep learning, and graph networks (2018). \texttt{arXiv:1806.01261}
  \item Hierarchical Graph Representation Learning with Differentiable Pooling (2018). \texttt{arXiv:1806.08804}
  \item Molecular Sets (MOSES): A Benchmarking Platform for Molecular Generation
  Models (2018). \texttt{arXiv:1811.12823}
  \item End-to-End Learning on 3D Protein Structure for Interface Prediction (2018). \texttt{arXiv:1807.01297}
  \item Deep convolutional networks for quality assessment of protein folds (2018). \texttt{arXiv:1801.06252}
  \item Graph Neural Networks: A Review of Methods and Applications (2018). \texttt{arXiv:1812.08434}
  \item NeVAE: A Deep Generative Model for Molecular Graphs (2018). \texttt{arXiv:1802.05283}
  \item Weisfeiler and Leman Go Neural: Higher-order Graph Neural Networks (2018). \texttt{arXiv:1810.02244}
  \item Representation Learning on Graphs with Jumping Knowledge Networks (2018). \texttt{arXiv:1806.03536}
  \item PADME: A Deep Learning-based Framework for Drug-Target Interaction
  Prediction (2018). \texttt{arXiv:1807.09741}
  \item Constrained Generation of Semantically Valid Graphs via Regularizing
  Variational Autoencoders (2018). \texttt{arXiv:1809.02630}
  \item Computational Protein Design with Deep Learning Neural Networks (2018). \texttt{arXiv:1801.07130}
  \item Deep Graph Infomax (2018). \texttt{arXiv:1809.10341}
  \item PotentialNet for Molecular Property Prediction (2018). \texttt{arXiv:1803.04465}
  \item Optimization of Molecules via Deep Reinforcement Learning (2018). \texttt{arXiv:1810.08678}
  \item Modeling polypharmacy side effects with graph convolutional networks (2018). \texttt{arXiv:1802.00543}
  \item Glow: Generative Flow with Invertible 1x1 Convolutions (2018). \texttt{arXiv:1807.03039}
  \item Learning Multimodal Graph-to-Graph Translation for Molecular
  Optimization (2018). \texttt{arXiv:1812.01070}
  \item Machine learning-guided directed evolution for protein engineering (2018). \texttt{arXiv:1811.10775}
  \item Graph Transformation Policy Network for Chemical Reaction Prediction (2018). \texttt{arXiv:1812.09441}
  \item PyTorch: An Imperative Style, High-Performance Deep Learning Library (2019). \texttt{arXiv:1912.01703}
  \item Symmetry-adapted generation of 3d point sets for the targeted discovery
  of molecules (2019). \texttt{arXiv:1906.00957}
  \item Fast Graph Representation Learning with PyTorch Geometric (2019). \texttt{arXiv:1903.02428}
  \item Analyzing Learned Molecular Representations for Property Prediction (2019). \texttt{arXiv:1904.01561}
  \item Efficient Graph Generation with Graph Recurrent Attention Networks (2019). \texttt{arXiv:1910.00760}
  \item Strategies for Pre-training Graph Neural Networks (2019). \texttt{arXiv:1905.12265}
  \item A Comprehensive Survey on Graph Neural Networks (2019). \texttt{arXiv:1901.00596}
  \item WideDTA: prediction of drug-target binding affinity (2019). \texttt{arXiv:1902.04166}
  \item Learning protein sequence embeddings using information from structure (2019). \texttt{arXiv:1902.08661}
  \item Molecular geometry prediction using a deep generative graph neural
  network (2019). \texttt{arXiv:1904.00314}
  \item GraphNVP: An Invertible Flow Model for Generating Molecular Graphs (2019). \texttt{arXiv:1905.11600}
  \item Evaluating Protein Transfer Learning with TAPE (2019). \texttt{arXiv:1906.08230}
  \item MolecularRNN: Generating realistic molecular graphs with optimized
  properties (2019). \texttt{arXiv:1905.13372}
  \item Optuna: A Next-generation Hyperparameter Optimization Framework (2019). \texttt{arXiv:1907.10902}
  \item A Two-Step Graph Convolutional Decoder for Molecule Generation (2019). \texttt{arXiv:1906.03412}
  \item ProteinNet: a standardized data set for machine learning of protein
  structure (2019). \texttt{arXiv:1902.00249}
  \item GNNExplainer: Generating Explanations for Graph Neural Networks (2019). \texttt{arXiv:1903.03894}
  \item InfoGraph: Unsupervised and Semi-supervised Graph-Level Representation
  Learning via Mutual Information Maximization (2019). \texttt{arXiv:1908.01000}
  \item Molecular Property Prediction: A Multilevel Quantum Interactions
  Modeling Perspective (2019). \texttt{arXiv:1906.11081}
  \item Self-Referencing Embedded Strings (SELFIES): A 100\% robust molecular
  string representation (2019). \texttt{arXiv:1905.13741}
  \item Graph Neural Network Architecture Search for Molecular Property
  Prediction (2020). \texttt{arXiv:2008.12187}
  \item Utilising Graph Machine Learning within Drug Discovery and Development (2020). \texttt{arXiv:2012.05716}
  \item Molecular graph generation with Graph Neural Networks (2020). \texttt{arXiv:2012.07397}
  \item Distance-aware Molecule Graph Attention Network for Drug-Target Binding
  Affinity Prediction (2020). \texttt{arXiv:2012.09624}
  \item Learn molecular representations from large-scale unlabeled molecules for
  drug discovery (2020). \texttt{arXiv:2012.11175}
  \item Learning from Protein Structure with Geometric Vector Perceptrons (2020). \texttt{arXiv:2009.01411}
  \item Directional Message Passing for Molecular Graphs (2020). \texttt{arXiv:2003.03123}
  \item GraphAF: a Flow-based Autoregressive Model for Molecular Graph
  Generation (2020). \texttt{arXiv:2001.09382}
  \item Language Models are Few-Shot Learners (2020). \texttt{arXiv:2005.14165}
  \item MoFlow: An Invertible Flow Model for Generating Molecular Graphs (2020). \texttt{arXiv:2006.10137}
  \item Symmetry-Aware Actor-Critic for 3D Molecular Design (2020). \texttt{arXiv:2011.12747}
  \item Permutation Invariant Graph Generation via Score-Based Generative
  Modeling (2020). \texttt{arXiv:2003.00638}
  \item Hierarchical Generation of Molecular Graphs using Structural Motifs (2020). \texttt{arXiv:2002.03230}
  \item ATOM3D: Tasks On Molecules in Three Dimensions (2020). \texttt{arXiv:2012.04035}
  \item Open Graph Benchmark: Datasets for Machine Learning on Graphs (2020). \texttt{arXiv:2005.00687}
  \item Molecule Attention Transformer (2020). \texttt{arXiv:2002.08264}
  \item Intrinsic-Extrinsic Convolution and Pooling for Learning on 3D Protein
  Structures (2020). \texttt{arXiv:2007.06252}
  \item Self-Supervised Graph Transformer on Large-Scale Molecular Data (2020). \texttt{arXiv:2007.02835}
  \item Graph Contrastive Learning with Augmentations (2020). \texttt{arXiv:2010.13902}
  \item Multi-Objective Molecule Generation using Interpretable Substructures (2020). \texttt{arXiv:2002.03244}
  \item Learning a Continuous Representation of 3D Molecular Structures with
  Deep Generative Models (2020). \texttt{arXiv:2010.08687}
  \item ProtTrans: Towards Cracking the Language of Life's Code Through
  Self-Supervised Deep Learning and High Performance Computing (2020). \texttt{arXiv:2007.06225}
  \item Deep Learning of High-Order Interactions for Protein Interface
  Prediction (2020). \texttt{arXiv:2007.09334}
  \item Relevance of Rotationally Equivariant Convolutions for Predicting
  Molecular Properties (2020). \texttt{arXiv:2008.08461}
  \item ChemBERTa: Large-Scale Self-Supervised Pretraining for Molecular
  Property Prediction (2020). \texttt{arXiv:2010.09885}
  \item Molecular Graph Generation via Geometric Scattering (2021). \texttt{arXiv:2110.06241}
  \item Quantitative Evaluation of Explainable Graph Neural Networks for
  Molecular Property Prediction (2021). \texttt{arXiv:2107.04119}
  \item Few-Shot Graph Learning for Molecular Property Prediction (2021). \texttt{arXiv:2102.07916}
  \item GeomGCL: Geometric Graph Contrastive Learning for Molecular Property
  Prediction (2021). \texttt{arXiv:2109.11730}
  \item Structure-aware Interactive Graph Neural Networks for the Prediction of
  Protein-Ligand Binding Affinity (2021). \texttt{arXiv:2107.10670}
  \item Learning Neural Generative Dynamics for Molecular Conformation
  Generation (2021). \texttt{arXiv:2102.10240}
  \item Pre-training Molecular Graph Representation with 3D Geometry (2021). \texttt{arXiv:2110.07728}
  \item Top-N: Equivariant set and graph generation without exchangeability (2021). \texttt{arXiv:2110.02096}
  \item Neural Algorithmic Reasoning (2021). \texttt{arXiv:2105.02761}
  \item Equivariant Graph Neural Networks for 3D Macromolecular Structure (2021). \texttt{arXiv:2106.03843}
  \item Molecular Contrastive Learning of Representations via Graph Neural
  Networks (2021). \texttt{arXiv:2102.10056}
  \item Independent SE(3)-Equivariant Models for End-to-End Rigid Protein
  Docking (2021). \texttt{arXiv:2111.07786}
  \item Therapeutics Data Commons: Machine Learning Datasets and Tasks for Drug
  Discovery and Development (2021). \texttt{arXiv:2102.09548}
  \item 3D Infomax improves GNNs for Molecular Property Prediction (2021). \texttt{arXiv:2110.04126}
  \item Motif-based Graph Self-Supervised Learning for Molecular Property
  Prediction (2021). \texttt{arXiv:2110.00987}
  \item Geometric Deep Learning on Molecular Representations (2021). \texttt{arXiv:2107.12375}
  \item GRAND: Graph Neural Diffusion (2021). \texttt{arXiv:2106.10934}
  \item Binary classification of proteins by a Machine Learning approach (2021). \texttt{arXiv:2111.01975}
  \item Application of Graph Neural Networks and graph descriptors for graph
  classification (2022). \texttt{arXiv:2211.03666}
  \item Graph neural networks for the prediction of molecular structure-property
  relationships (2022). \texttt{arXiv:2208.04852}
  \item Graph Neural Networks for Molecules (2022). \texttt{arXiv:2209.05582}
  \item HiGNN: Hierarchical Informative Graph Neural Networks for Molecular
  Property Prediction Equipped with Feature-Wise Attention (2022). \texttt{arXiv:2208.13994}
  \item Graph neural networks for materials science and chemistry (2022). \texttt{arXiv:2208.09481}
  \item Equivariant Diffusion for Molecule Generation in 3D (2022). \texttt{arXiv:2203.17003}
  \item GeoDiff: a Geometric Diffusion Model for Molecular Conformation
  Generation (2022). \texttt{arXiv:2203.02923}
  \item Torsional Diffusion for Molecular Conformer Generation (2022). \texttt{arXiv:2206.01729}
  \item A 3D Generative Model for Structure-Based Drug Design (2022). \texttt{arXiv:2203.10446}
  \item DiGress: Discrete Denoising diffusion for graph generation (2022). \texttt{arXiv:2209.14734}
  \item Structure-based Drug Design with Equivariant Diffusion Models (2022). \texttt{arXiv:2210.13695}
  \item Pocket2Mol: Efficient Molecular Sampling Based on 3D Protein Pockets (2022). \texttt{arXiv:2205.07249}
  \item DiffDock: Diffusion Steps, Twists, and Turns for Molecular Docking (2022). \texttt{arXiv:2210.01776}
  \item Generating 3D Molecules for Target Protein Binding (2022). \texttt{arXiv:2204.09410}
  \item Protein Representation Learning by Geometric Structure Pretraining (2022). \texttt{arXiv:2203.06125}
  \item Diffusion probabilistic modeling of protein backbones in 3D for the
  motif-scaffolding problem (2022). \texttt{arXiv:2206.04119}
  \item Learning Hierarchical Protein Representations via Complete 3D Graph
  Networks (2022). \texttt{arXiv:2207.12600}
  \item EquiBind: Geometric Deep Learning for Drug Binding Structure Prediction (2022). \texttt{arXiv:2202.05146}
  \item Protein Structure and Sequence Generation with Equivariant Denoising
  Diffusion Probabilistic Models (2022). \texttt{arXiv:2205.15019}
  \item Diffusion-based Molecule Generation with Informative Prior Bridges (2022). \texttt{arXiv:2209.00865}
  \item Equivariant 3D-Conditional Diffusion Models for Molecular Linker Design (2022). \texttt{arXiv:2210.05274}
  \item Multi-Scale Representation Learning on Proteins (2022). \texttt{arXiv:2204.02337}
  \item DiffBP: Generative Diffusion of 3D Molecules for Target Protein Binding (2022). \texttt{arXiv:2211.11214}
  \item SPECTRE: Spectral Conditioning Helps to Overcome the Expressivity Limits
  of One-shot Graph Generators (2022). \texttt{arXiv:2204.01613}
  \item Diffusion Models: A Comprehensive Survey of Methods and Applications (2022). \texttt{arXiv:2209.00796}
  \item Protein structure generation via folding diffusion (2022). \texttt{arXiv:2209.15611}
  \item Hierarchical Graph Representation Learning for the Prediction of
  Drug-Target Binding Affinity (2022). \texttt{arXiv:2203.11458}
  \item MDM: Molecular Diffusion Model for 3D Molecule Generation (2022). \texttt{arXiv:2209.05710}
  \item A Generalist Neural Algorithmic Learner (2022). \texttt{arXiv:2209.11142}
  \item Equivariant Energy-Guided SDE for Inverse Molecular Design (2022). \texttt{arXiv:2209.15408}
  \item Direct Molecular Conformation Generation (2022). \texttt{arXiv:2202.01356}
  \item Translation between Molecules and Natural Language (2022). \texttt{arXiv:2204.11817}
  \item Graph Neural Networks are Dynamic Programmers (2022). \texttt{arXiv:2203.15544}
  \item ComENet: Towards Complete and Efficient Message Passing for 3D Molecular
  Graphs (2022). \texttt{arXiv:2206.08515}
  \item 3DLinker: An E(3) Equivariant Variational Autoencoder for Molecular
  Linker Design (2022). \texttt{arXiv:2205.07309}
  \item Molecular Property Prediction Based on Graph Structure Learning (2023). \texttt{arXiv:2312.16855}
  \item Building explainable graph neural network by sparse learning for the
  drug-protein binding prediction (2023). \texttt{arXiv:2309.12906}
  \item A Survey on Graph Diffusion Models: Generative AI in Science for
  Molecule, Protein and Material (2023). \texttt{arXiv:2304.01565}
  \item Knowledge-augmented Graph Machine Learning for Drug Discovery: A Survey (2023). \texttt{arXiv:2302.08261}
  \item DrugChat: Towards Enabling ChatGPT-Like Capabilities on Drug Molecule
  Graphs (2023). \texttt{arXiv:2309.03907}
  \item Geometric Latent Diffusion Models for 3D Molecule Generation (2023). \texttt{arXiv:2305.01140}
  \item MiDi: Mixed Graph and 3D Denoising Diffusion for Molecule Generation (2023). \texttt{arXiv:2302.09048}
  \item 3D Equivariant Diffusion for Target-Aware Molecule Generation and
  Affinity Prediction (2023). \texttt{arXiv:2303.03543}
  \item MolDiff: Addressing the Atom-Bond Inconsistency Problem in 3D Molecule
  Diffusion Generation (2023). \texttt{arXiv:2305.07508}
  \item SE(3) diffusion model with application to protein backbone generation (2023). \texttt{arXiv:2302.02277}
  \item LLaMA: Open and Efficient Foundation Language Models (2023). \texttt{arXiv:2302.13971}
  \item Autoregressive Diffusion Model for Graph Generation (2023). \texttt{arXiv:2307.08849}
  \item GraphGANFed: A Federated Generative Framework for Graph-Structured
  Molecules Towards Efficient Drug Discovery (2023). \texttt{arXiv:2304.05498}
  \item MUDiff: Unified Diffusion for Complete Molecule Generation (2023). \texttt{arXiv:2304.14621}
  \item Learning Joint 2D \& 3D Diffusion Models for Complete Molecule Generation (2023). \texttt{arXiv:2305.12347}
  \item Using GNN property predictors as molecule generators (2024). \texttt{arXiv:2406.03278}
  \item On the Scalability of GNNs for Molecular Graphs (2024). \texttt{arXiv:2404.11568}
  \item Molecular Graph Representation Learning via Structural Similarity
  Information (2024). \texttt{arXiv:2409.08580}
  \item Graph Neural Networks for Protein-Protein Interactions -- A Short Survey (2024). \texttt{arXiv:2404.10450}
  \item ViDTA: Enhanced Drug-Target Affinity Prediction via Virtual Graph Nodes
  and Attention-based Feature Fusion (2024). \texttt{arXiv:2412.19589}
  \item DecompDiff: Diffusion Models with Decomposed Priors for Structure-Based
  Drug Design (2024). \texttt{arXiv:2403.07902}
  \item Tutorial on Diffusion Models for Imaging and Vision (2024). \texttt{arXiv:2403.18103}
  \item Closed-form estimators for an exponential family derived from likelihood
  equations (2024). \texttt{arXiv:2405.14509}
  \item Graph Neural Networks in Modern AI-aided Drug Discovery (2025). \texttt{arXiv:2506.06915}
  \item A Survey of Graph Neural Networks for Drug Discovery: Recent Developments and Challenges (2025). \texttt{arXiv:2509.07887}
  \item Recent Developments in GNNs for Drug Discovery (2025). \texttt{arXiv:2506.01302}
  \item Multi-Level Fusion Graph Neural Network for Molecule Property Prediction (2025). \texttt{arXiv:2507.03430}
  \item A Graph-in-Graph Learning Framework for Drug-Target Interaction Prediction (2025). \texttt{arXiv:2507.11757}
  \item Geometry and biological processes (2026). \texttt{arXiv:2602.20274}
  \item Invariants in equivariant birational geometry (2026). \texttt{arXiv:2602.23998}
\end{enumerate}

\end{document}
